\documentclass[12pt]{article}
\usepackage[utf8]{inputenc}
\usepackage[T1]{fontenc}
\usepackage[italian]{babel}
\usepackage{amsmath, amssymb, amsthm}
\usepackage{geometry}
\usepackage{booktabs}
\usepackage{array}
\usepackage{enumitem}
\usepackage{bm}
\usepackage{graphicx}
\usepackage{xcolor}
\usepackage{hyperref}

\geometry{a4paper, left=25mm, right=25mm, top=30mm, bottom=30mm}

\title{\textbf{Spline Quintiche di Hermite: \\ Calcolo delle Derivate Seconde per Continuità $C^3$}}
\author{Metodi Numerici per l'Interpolazione}
\date{\today}

\newcommand{\R}{\mathbb{R}}
\newcommand{\pd}{\partial}
\newcommand{\pder}[2]{\frac{\pd #1}{\pd #2}}

\begin{document}

\maketitle

\begin{abstract}
Questo documento descrive il procedimento per calcolare le derivate seconde nei nodi di una spline composta da polinomi di Hermite di grado 5, con le derivate prime fissate, in modo da garantire la continuità delle derivate terze. Viene presentata la derivazione matematica completa, le diverse condizioni al contorno disponibili e l'algoritmo numerico per la risoluzione del sistema lineare tridiagonale risultante.
\end{abstract}

\tableofcontents

\section{Introduzione}
Le spline di Hermite di grado 5 sono uno strumento potente per l'interpolazione e l'approssimazione di funzioni, offrendo continuità $C^2$ di default e la possibilità di ottenere continuità $C^3$ con un'opportuna scelta delle derivate seconde. Questo documento affronta il problema specifico di determinare le derivate seconde nei nodi quando le derivate prime sono già fissate, imponendo la continuità delle derivate terze.

\section{Formulazione del Problema}

\subsection{Dati iniziali}
Si considerino $n+1$ nodi distinti:
\[
x_0 < x_1 < \dots < x_n, \quad x_i \in \R
\]
con valori della funzione e derivate prime note:
\[
f_i = f(x_i), \quad d_i = f'(x_i), \quad i = 0,\dots,n.
\]
Si definiscano le lunghezze degli intervalli:
\[
h_i = x_{i+1} - x_i > 0, \quad i = 0,\dots,n-1.
\]

\subsection{Polinomio di Hermite di grado 5}
In ogni intervallo $[x_i, x_{i+1}]$, il polinomio $p_i(x)$ di grado 5 è univocamente determinato dai sei valori:
\[
p_i(x_i) = f_i,\quad p_i(x_{i+1}) = f_{i+1},
\]
\[
p_i'(x_i) = d_i,\quad p_i'(x_{i+1}) = d_{i+1},
\]
\[
p_i''(x_i) = s_i,\quad p_i''(x_{i+1}) = s_{i+1},
\]
dove $s_i = f''(x_i)$ sono le incognite da determinare.

\section{Derivazione delle Equazioni di Continuità}

\subsection{Rappresentazione normalizzata}
Introduciamo il parametro normalizzato $t \in [0,1]$:
\[
t = \frac{x - x_i}{h_i}, \quad x \in [x_i, x_{i+1}].
\]
Il polinomio $p_i(t)$ può essere espresso come combinazione lineare dei polinomi base di Hermite:
\[
\begin{aligned}
p_i(t) &= f_i H_0(t) + f_{i+1} H_1(t) + h_i d_i H_2(t) + h_i d_{i+1} H_3(t) \\
&\quad + h_i^2 s_i H_4(t) + h_i^2 s_{i+1} H_5(t),
\end{aligned}
\]
con:
\[
\begin{aligned}
H_0(t) &= 1 - 10t^3 + 15t^4 - 6t^5, \\
H_1(t) &= 10t^3 - 15t^4 + 6t^5, \\
H_2(t) &= t - 6t^3 + 8t^4 - 3t^5, \\
H_3(t) &= -4t^3 + 7t^4 - 3t^5, \\
H_4(t) &= \frac{1}{2}(t^2 - 3t^3 + 3t^4 - t^5), \\
H_5(t) &= \frac{1}{2}(-t^3 + 2t^4 - t^5).
\end{aligned}
\]

\subsection{Derivate terze}
La derivata terza rispetto a $x$ è:
\[
p_i'''(x) = \frac{1}{h_i^3} p_i'''(t).
\]
Calcoliamo le derivate terze dei polinomi base:
\[
\begin{aligned}
H_0'''(t) &= -60 + 180t - 120t^2, \\
H_1'''(t) &= 60 - 180t + 120t^2, \\
H_2'''(t) &= -36 + 96t - 60t^2, \\
H_3'''(t) &= -24 + 84t - 60t^2, \\
H_4'''(t) &= \frac{1}{2}(-18 + 72t - 60t^2) = -9 + 36t - 30t^2, \\
H_5'''(t) &= \frac{1}{2}(-6 + 24t - 20t^2) = -3 + 12t - 10t^2.
\end{aligned}
\]

\subsection{Condizioni di continuità $C^3$}
Nei nodi interni $x_i$, $i=1,\dots,n-1$, imponiamo:
\[
p_{i-1}'''(x_i^-) = p_i'''(x_i^+).
\]

\subsubsection{Valutazione a sinistra ($t=1$ per $p_{i-1}$)}
\[
\begin{aligned}
p_{i-1}'''(1) &= f_{i-1} H_0'''(1) + f_i H_1'''(1) + h_{i-1} d_{i-1} H_2'''(1) + h_{i-1} d_i H_3'''(1) \\
&\quad + h_{i-1}^2 s_{i-1} H_4'''(1) + h_{i-1}^2 s_i H_5'''(1).
\end{aligned}
\]
Valutando:
\[
\begin{aligned}
H_0'''(1) &= 0, \quad H_1'''(1) = 0, \quad H_2'''(1) = 0, \quad H_3'''(1) = 0, \\
H_4'''(1) &= -3, \quad H_5'''(1) &= -1.
\end{aligned}
\]
Quindi:
\[
p_{i-1}'''(x_i^-) = \frac{1}{h_{i-1}^3} \left(-3 h_{i-1}^2 s_{i-1} - h_{i-1}^2 s_i\right) = -\frac{1}{h_{i-1}}(3s_{i-1} + s_i).
\]

\subsubsection{Valutazione a destra ($t=0$ per $p_i$)}
\[
\begin{aligned}
p_i'''(0) &= f_i H_0'''(0) + f_{i+1} H_1'''(0) + h_i d_i H_2'''(0) + h_i d_{i+1} H_3'''(0) \\
&\quad + h_i^2 s_i H_4'''(0) + h_i^2 s_{i+1} H_5'''(0).
\end{aligned}
\]
Valutando:
\[
\begin{aligned}
H_0'''(0) &= -60, \quad H_1'''(0) = 60, \\
H_2'''(0) &= -36, \quad H_3'''(0) = -24, \\
H_4'''(0) &= -9, \quad H_5'''(0) = -3.
\end{aligned}
\]
Quindi:
\[
p_i'''(x_i^+) = \frac{1}{h_i^3} \left[ -60f_i + 60f_{i+1} - 36h_i d_i - 24h_i d_{i+1} - 9h_i^2 s_i - 3h_i^2 s_{i+1} \right].
\]

\subsection{Equazione risultante}
Uguagliando $p_{i-1}'''(x_i^-) = p_i'''(x_i^+)$ e moltiplicando per $h_{i-1}^3 h_i^3$:
\[
\begin{aligned}
&-h_i^3(3s_{i-1} + s_i) = h_{i-1}^3 \left[ -60f_i + 60f_{i+1} - 36h_i d_i - 24h_i d_{i+1} - 9h_i^2 s_i - 3h_i^2 s_{i+1} \right].
\end{aligned}
\]
Riorganizzando:
\[
- h_i^3 s_{i-1} + 3h_{i-1}^3 h_i s_i - h_{i-1}^3 s_{i+1} = R_i,
\]
dove:
\[
\begin{aligned}
R_i = -60h_{i-1}^3(f_{i+1} - f_i) - 36h_{i-1}^3 h_i d_i - 24h_{i-1}^3 h_i d_{i+1}.
\end{aligned}
\]

\subsection{Forma standard}
Un'ulteriore manipolazione algebrica porta alla forma standard tridiagonale:
\[
\boxed{- h_i s_{i-1} + 3(h_i + h_{i-1}) s_i - h_{i-1} s_{i+1} = \tilde{R}_i},
\]
per $i = 1,\dots,n-1$, con:
\[
\boxed{\tilde{R}_i = -20\left[ \frac{h_i(f_i - f_{i-1})}{h_{i-1}^2} - \frac{h_{i-1}(f_{i+1} - f_i)}{h_i^2} \right] - 12 \frac{h_i^2 + h_{i-1}^2}{h_{i-1} h_i} d_i - 8\left( \frac{h_i}{h_{i-1}} d_{i-1} + \frac{h_{i-1}}{h_i} d_{i+1} \right)}.
\]

\section{Condizioni al Contorno}

Il sistema fornisce $n-1$ equazioni per $n+1$ incognite $s_0,\dots,s_n$. Sono necessarie due condizioni aggiuntive.

\subsection{Tipo I: Derivate seconde specificate}
Assegnamo valori noti:
\[
s_0 = \alpha, \quad s_n = \beta.
\]
Il sistema diventa:
\[
\begin{bmatrix}
B_1 & C_1 & 0 & \cdots & 0 \\
A_2 & B_2 & C_2 & \cdots & 0 \\
0 & \ddots & \ddots & \ddots & 0 \\
0 & \cdots & A_{n-2} & B_{n-2} & C_{n-2} \\
0 & \cdots & 0 & A_{n-1} & B_{n-1}
\end{bmatrix}
\begin{bmatrix}
s_1 \\ s_2 \\ \vdots \\ s_{n-2} \\ s_{n-1}
\end{bmatrix}
=
\begin{bmatrix}
\tilde{R}_1 - A_1\alpha \\ \tilde{R}_2 \\ \vdots \\ \tilde{R}_{n-2} \\ \tilde{R}_{n-1} - C_{n-1}\beta
\end{bmatrix},
\]
dove:
\[
A_i = -h_i, \quad B_i = 3(h_i + h_{i-1}), \quad C_i = -h_{i-1}.
\]

\subsection{Tipo II: Spline naturale (derivate seconde nulle)}
Caso particolare del Tipo I con:
\[
s_0 = 0, \quad s_n = 0.
\]

\subsection{Tipo III: Derivate terze nulle agli estremi}
Imponiamo $p_0'''(x_0) = 0$ e $p_{n-1}'''(x_n) = 0$.

\subsubsection{Estremo sinistro}
Da $p_0'''(0) = 0$:
\[
2h_0 s_0 + h_0 s_1 = -10\frac{f_1 - f_0}{h_0} + 6d_0 + 4d_1.
\]

\subsubsection{Estremo destro}
Da $p_{n-1}'''(1) = 0$:
\[
h_{n-1} s_{n-1} + 2h_{n-1} s_n = 10\frac{f_n - f_{n-1}}{h_{n-1}} - 4d_{n-1} - 6d_n.
\]

Il sistema completo ($n+1$ equazioni in $n+1$ incognite) è:
\[
\begin{bmatrix}
2h_0 & h_0 & 0 & \cdots & 0 \\
A_1 & B_1 & C_1 & \cdots & 0 \\
0 & \ddots & \ddots & \ddots & 0 \\
0 & \cdots & A_{n-1} & B_{n-1} & C_{n-1} \\
0 & \cdots & 0 & h_{n-1} & 2h_{n-1}
\end{bmatrix}
\begin{bmatrix}
s_0 \\ s_1 \\ \vdots \\ s_{n-1} \\ s_n
\end{bmatrix}
=
\begin{bmatrix}
Q_0 \\ \tilde{R}_1 \\ \vdots \\ \tilde{R}_{n-1} \\ Q_n
\end{bmatrix},
\]
con:
\[
Q_0 = -10\frac{f_1 - f_0}{h_0} + 6d_0 + 4d_1, \quad Q_n = 10\frac{f_n - f_{n-1}}{h_{n-1}} - 4d_{n-1} - 6d_n.
\]

\subsection{Tipo IV: Condizioni periodiche}
Se $f_0 = f_n$, $d_0 = d_n$, e $s_0 = s_n$, il sistema diventa ciclico. Le equazioni per $i=1$ e $i=n$ diventano:
\[
\begin{aligned}
- h_1 s_n + 3(h_1 + h_0) s_1 - h_0 s_2 &= \tilde{R}_1, \\
- h_0 s_{n-1} + 3(h_0 + h_{n-1}) s_n - h_{n-1} s_1 &= \tilde{R}_n,
\end{aligned}
\]
dove si è posto $h_n = h_0$ e $\tilde{R}_n$ si calcola con le opportune sostituzioni.

\section{Algoritmo di Soluzione}

\subsection{Sistema tridiagonale}
Il sistema per le condizioni Tipo I/II ha matrice tridiagonale strettamente diagonalmente dominante per $h_i > 0$, garantendo un'unica soluzione.

\subsection{Algoritmo di Thomas (TDMA)}
Per risolvere il sistema $A\bm{s} = \bm{D}$ con:
\[
A = \begin{bmatrix}
B_1 & C_1 & & & \\
A_2 & B_2 & C_2 & & \\
& \ddots & \ddots & \ddots & \\
& & A_{m-1} & B_{m-1} & C_{m-1} \\
& & & A_m & B_m
\end{bmatrix}, \quad \bm{s} = \begin{bmatrix} s_1 \\ s_2 \\ \vdots \\ s_{m-1} \\ s_m \end{bmatrix}, \quad \bm{D} = \begin{bmatrix} D_1 \\ D_2 \\ \vdots \\ D_{m-1} \\ D_m \end{bmatrix}.
\]

\subsubsection{Forward sweep}
Calcolare:
\[
\begin{aligned}
c'_1 &= \frac{C_1}{B_1}, \quad d'_1 = \frac{D_1}{B_1}, \\
c'_i &= \frac{C_i}{B_i - A_i c'_{i-1}}, \quad i = 2,\dots,m-1, \\
d'_i &= \frac{D_i - A_i d'_{i-1}}{B_i - A_i c'_{i-1}}, \quad i = 2,\dots,m.
\end{aligned}
\]

\subsubsection{Backward substitution}
\[
\begin{aligned}
s_m &= d'_m, \\
s_i &= d'_i - c'_i s_{i+1}, \quad i = m-1,\dots,1.
\end{aligned}
\]

\subsection{Pseudocodice per condizioni Tipo I/II}
\begin{verbatim}
Input: x[0..n], f[0..n], d[0..n], alpha, beta
Output: s[0..n]

// Calcolo lunghezze intervalli
for i = 0 to n-1:
    h[i] = x[i+1] - x[i]

// Costruzione vettore dei termini noti D[1..n-1]
for i = 1 to n-1:
    term1 = -20 * ( h[i]*(f[i]-f[i-1])/(h[i-1]^2) 
                    - h[i-1]*(f[i+1]-f[i])/(h[i]^2) )
    term2 = -12 * (h[i]^2 + h[i-1]^2)/(h[i-1]*h[i]) * d[i]
    term3 = -8 * ( (h[i]/h[i-1])*d[i-1] + (h[i-1]/h[i])*d[i+1] )
    D[i] = term1 + term2 + term3

// Aggiusta per condizioni al contorno
D[1] = D[1] + h[1]*alpha
D[n-1] = D[n-1] + h[n-2]*beta

// Coefficienti della matrice
for i = 1 to n-1:
    A[i] = -h[i]           // sotto-diagonale (indici: i, i-1)
    B[i] = 3*(h[i] + h[i-1]) // diagonale
    C[i] = -h[i-1]         // sopra-diagonale (indici: i, i+1)

// Risoluzione sistema tridiagonale con Thomas
// ... (implementare algoritmo sopra)

// Assegna valori noti
s[0] = alpha
s[n] = beta
\end{verbatim}

\section{Proprietà della Soluzione}

\subsection{Esistenza e unicità}
\begin{itemize}
\item \textbf{Tipo I/II}: La matrice è diagonalmente dominante per $h_i > 0$, garantendo un'unica soluzione.
\item \textbf{Tipo III}: Condizioni simili garantiscono l'unicità per $n \geq 2$.
\item \textbf{Tipo IV}: Richiede che il sistema ciclico sia non singolare (condizione tipicamente soddisfatta).
\end{itemize}

\subsection{Stabilità numerica}
L'algoritmo di Thomas è stabile per matrici diagonalmente dominanti. In caso di griglie uniformi ($h_i = h$), la matrice ha coefficienti costanti, migliorando ulteriormente la stabilità.

\subsection{Complessità computazionale}
\begin{itemize}
\item Calcolo dei coefficienti: $O(n)$.
\item Algoritmo di Thomas: $O(n)$ operazioni.
\item Complessità totale: $O(n)$.
\end{itemize}

\section{Esempio Numerico}

\subsection{Dati}
Consideriamo nodi $x = [0, 1, 2, 3]$, valori $f = [0, 1, 0, 1]$, derivate prime $d = [1, 0, -1, 0]$, e condizioni naturali $s_0 = s_3 = 0$.

\subsection{Calcolo}
Con $h_0 = h_1 = h_2 = 1$, il sistema per $s_1, s_2$ è:
\[
\begin{bmatrix}
6 & -1 \\ -1 & 6
\end{bmatrix}
\begin{bmatrix}
s_1 \\ s_2
\end{bmatrix}
=
\begin{bmatrix}
-20(1-0) - 20(0-1) - 12(1+1)\cdot 0 - 8(1\cdot 1 + 1\cdot (-1)) \\
-20(0-1) - 20(1-0) - 12(1+1)\cdot (-1) - 8(1\cdot 0 + 1\cdot 0)
\end{bmatrix}
=
\begin{bmatrix}
0 \\ 24
\end{bmatrix}.
\]
Soluzione: $s_1 = \frac{24}{35} \approx 0.6857$, $s_2 = \frac{144}{35} \approx 4.1143$.

\section{Conclusioni}
Il metodo descritto permette di calcolare efficientemente le derivate seconde per spline quintiche di Hermite con continuità $C^3$. La scelta delle condizioni al contorno dipende dall'applicazione:
\begin{itemize}
\item \textbf{Tipo I}: Quando si conoscono le curvature agli estremi.
\item \textbf{Tipo II}: Spline naturale, semplice ma può introdurre oscillazioni.
\item \textbf{Tipo III}: Comportamento "liscio" agli estremi, spesso preferibile.
\item \textbf{Tipo IV}: Per dati periodici.
\end{itemize}
L'implementazione con algoritmo di Thomas garantisce efficienza e stabilità numerica.

\section*{Riferimenti}
\begin{enumerate}
\item Carl de Boor, \emph{A Practical Guide to Splines}, Springer-Verlag, 1978.
\item David Kincaid, Ward Cheney, \emph{Numerical Analysis}, Brooks/Cole, 2002.
\item John H. Mathews, Kurtis D. Fink, \emph{Numerical Methods Using MATLAB}, Prentice Hall, 2004.
\end{enumerate}

\end{document}